% --- Archon physics formalization source begin ---
% archon:physics
% archon:covers PhyXMiniProblems/problem_phyx_mini_0562.lean
% archon:source-report reports/phyx_mini/problem_phyx_mini_0562.source.json

\chapter{Physics problem phyx\_mini\_0562}
\label{ch:PhyXMiniProblems_problem_phyx_mini_0562}

\paragraph{Problem source.}
\#\# Physical scenario

A small shot of negligible radius hits a stationary smooth, hard sphere of radius \textbackslash{}( R \textbackslash{}), making an angle \textbackslash{}( \textbackslash{}beta \textbackslash{}) with the normal to the sphere, as shown in figure. It is reflected at an equal angle to the normal. The scattering angle is \textbackslash{}( \textbackslash{}theta = 180\textasciicircum{}\{\textbackslash{}circ\} - 2\textbackslash{}beta \textbackslash{}), as shown.

\#\# Question

If the incoming intensity of the shot is \textbackslash{}( I\_0 \textbackslash{}) particles/s \textbackslash{}( \textbackslash{}cdot \textbackslash{}) area, how many are scattered through angles greater than \textbackslash{}( \textbackslash{}theta \textbackslash{})?

\#\# Answer choices

- A: A: \textbackslash{}( I\_0 R\textasciicircum{}2 \textbackslash{}sin\textasciicircum{}2 \textbackslash{}frac\{\textbackslash{}theta\}\{2\} \textbackslash{})

- B: B: \textbackslash{}( I\_0 R\textasciicircum{}2 \textbackslash{}tan\textasciicircum{}2 \textbackslash{}frac\{\textbackslash{}theta\}\{2\} \textbackslash{})

- C: C: \textbackslash{}( I\_0 R\textasciicircum{}2 \textbackslash{}cot\textasciicircum{}2 \textbackslash{}frac\{\textbackslash{}theta\}\{2\} \textbackslash{})

- D: D: \textbackslash{}( I\_0 R\textasciicircum{}2 \textbackslash{}cos\textasciicircum{}2 \textbackslash{}frac\{\textbackslash{}theta\}\{2\} \textbackslash{})

\#\# Figure caption (auxiliary; use the image as primary evidence)

The image is a geometric diagram involving a circle and several lines with angles and dimensions marked. Here's a breakdown of the contents:

- **Circle:** A blue circle is present on the right side, with a radius labeled as \textbackslash{}( R \textbackslash{}).
- **Horizontal Line:** There is a horizontal line passing tangentially from the center of the circle through to the left, labeled with a downward pointing arrow and distance \textbackslash{}( b \textbackslash{}) marked from the top of this line to the arrowhead of an upward vertical line.
- **Rays and Angles:**
  - A ray extends to the right from the horizontal line, forming an angle \textbackslash{}(\textbackslash{}beta\textbackslash{}) with a dashed line connecting a point on the ray to the center of the circle.
  - Another ray projects upwards from the tangent point on the circle, forming an angle \textbackslash{}(\textbackslash{}beta\textbackslash{}) with the dashed line.
  - An angle \textbackslash{}(\textbackslash{}theta\textbackslash{}) is marked between the upward ray and the horizontal line.
- **Points:** Three points are marked along the blue lines, two on the horizontal line and one on the extending ray.

The diagram seems to illustrate concepts related to geometry or physics, possibly dealing with angles of reflection or refraction.

\#\# Dataset metadata

Quantum Phenomena; Spatial Relation Reasoning, Physical Model Grounding Reasoning

\paragraph{Recorded answer/context.}
D: D: \textbackslash{}( I\_0 R\textasciicircum{}2 \textbackslash{}cos\textasciicircum{}2 \textbackslash{}frac\{\textbackslash{}theta\}\{2\} \textbackslash{})

\paragraph{Figure/image path.}
/root/proposal\_for\_physic/science-mango/phyx\_mini\_run/phyx\_data/test\_image/562.png

\paragraph{Formalization target.}
create a compiling Lean file with sorry bodies at `PhyXMiniProblems/problem\_phyx\_mini\_0562.lean`. The Lean declarations must preserve the physical quantities, dimensions or dimensional roles, figure labels, governing-law hypotheses, and final relation expressed by this problem.
Use Mathlib/Physlib names found through LeanExplore where available. If a physics API is missing, introduce faithful local abstractions rather than scalar placeholder aliases.

\begin{theorem}[Physics formalization target]
\label{thm:physics:phyx_mini_0562:target}
\lean{PhyXMiniProblems.ProblemPhyXMini0562.number_scattered_per_second_above_theta}
\uses{def:physics:phyx-mini-0562:phyxminiproblems-problemphyxmini0562-crosssectionnormalizationfactor, def:physics:phyx-mini-0562:phyxminiproblems-problemphyxmini0562-lengthinmeters, def:physics:phyx-mini-0562:phyxminiproblems-problemphyxmini0562-intensityinparticlespersecondsquaremeter, def:physics:phyx-mini-0562:phyxminiproblems-problemphyxmini0562-rateinparticlespersecond, def:physics:phyx-mini-0562:phyxminiproblems-problemphyxmini0562-hardspherescatteringsetup, def:physics:phyx-mini-0562:phyxminiproblems-problemphyxmini0562-matcheshardspherescatteringscenario, def:physics:phyx-mini-0562:phyxminiproblems-problemphyxmini0562-matchesprimaryfigure, def:physics:phyx-mini-0562:phyxminiproblems-problemphyxmini0562-hasphysicalhardsphereparameters, def:physics:phyx-mini-0562:phyxminiproblems-problemphyxmini0562-satisfiesspecularreflectionlaw, def:physics:phyx-mini-0562:phyxminiproblems-problemphyxmini0562-satisfieshardspherecontactgeometry, def:physics:phyx-mini-0562:phyxminiproblems-problemphyxmini0562-satisfiesuniformbeamcountinglaw, def:physics:phyx-mini-0562:phyxminiproblems-problemphyxmini0562-satisfieshardspherecrosssectionlaw}
The assigned autoformalize agent should translate this physics problem into Lean declarations in the covered file, with theorem and lemma proof bodies written as `by sorry`.
\end{theorem}
\begin{proof}
This is an autoformalization task, not a proof task. Produce statements that can later be proved by the physics prover without weakening the physical model.
\end{proof}

% --- Archon declaration topology begin ---
\section{Declaration topology}
\begin{definition}[Length Magnitude]
  \label{def:physics:phyx-mini-0562:phyxminiproblems-problemphyxmini0562-lengthmagnitude}
  \lean{PhyXMiniProblems.ProblemPhyXMini0562.LengthMagnitude}
  A nonnegative physical length, independent of the chosen length unit
\end{definition}
\begin{proof}
  This is the declared model definition of Length Magnitude.
\end{proof}
\begin{definition}[Area Magnitude]
  \label{def:physics:phyx-mini-0562:phyxminiproblems-problemphyxmini0562-areamagnitude}
  \lean{PhyXMiniProblems.ProblemPhyXMini0562.AreaMagnitude}
  A physical area, such as an effective scattering cross-section
\end{definition}
\begin{proof}
  This is the declared model definition of Area Magnitude.
\end{proof}
\begin{definition}[Particle Intensity]
  \label{def:physics:phyx-mini-0562:phyxminiproblems-problemphyxmini0562-particleintensity}
  \lean{PhyXMiniProblems.ProblemPhyXMini0562.ParticleIntensity}
  Incident particle intensity, with dimension particles per unit area per unit time. Particle number is dimensionless, so the physical dimension is (length² · time)⁻¹
\end{definition}
\begin{proof}
  This is the declared model definition of Particle Intensity.
\end{proof}
\begin{definition}[Particle Rate]
  \label{def:physics:phyx-mini-0562:phyxminiproblems-problemphyxmini0562-particlerate}
  \lean{PhyXMiniProblems.ProblemPhyXMini0562.ParticleRate}
  A number of particles per unit time
\end{definition}
\begin{proof}
  This is the declared model definition of Particle Rate.
\end{proof}
\begin{definition}[Cross Section Normalization]
  \label{def:physics:phyx-mini-0562:phyxminiproblems-problemphyxmini0562-crosssectionnormalization}
  \lean{PhyXMiniProblems.ProblemPhyXMini0562.CrossSectionNormalization}
  There are two normalization conventions relevant to the source. An ordinary three-dimensional beam counts the full impact disk and therefore carries a factor of π. Every displayed answer omits that factor, so the source can only be read literally after adopting its reduced (azimuth-normalized) convention. Keeping this choice as data prevents the ambiguity from being silently folded into the physical laws or the final answer
\end{definition}
\begin{proof}
  This is the declared model definition of Cross Section Normalization.
\end{proof}
\begin{definition}[cross Section Normalization Factor]
  \label{def:physics:phyx-mini-0562:phyxminiproblems-problemphyxmini0562-crosssectionnormalizationfactor}
  \lean{PhyXMiniProblems.ProblemPhyXMini0562.crossSectionNormalizationFactor}
  \uses{def:physics:phyx-mini-0562:phyxminiproblems-problemphyxmini0562-crosssectionnormalization}
  Dimensionless area factor belonging to a cross-section convention
\end{definition}
\begin{proof}
  This is the declared model definition of cross Section Normalization Factor.
\end{proof}
\begin{definition}[length Readout]
  \label{def:physics:phyx-mini-0562:phyxminiproblems-problemphyxmini0562-lengthreadout}
  \lean{PhyXMiniProblems.ProblemPhyXMini0562.lengthReadout}
  \uses{def:physics:phyx-mini-0562:phyxminiproblems-problemphyxmini0562-lengthmagnitude}
  Read a physical length in a selected length unit
\end{definition}
\begin{proof}
  This is the declared model definition of length Readout.
\end{proof}
\begin{definition}[area Readout]
  \label{def:physics:phyx-mini-0562:phyxminiproblems-problemphyxmini0562-areareadout}
  \lean{PhyXMiniProblems.ProblemPhyXMini0562.areaReadout}
  \uses{def:physics:phyx-mini-0562:phyxminiproblems-problemphyxmini0562-areamagnitude}
  Read a physical area in the square of a selected length unit
\end{definition}
\begin{proof}
  This is the declared model definition of area Readout.
\end{proof}
\begin{definition}[particle Intensity Readout]
  \label{def:physics:phyx-mini-0562:phyxminiproblems-problemphyxmini0562-particleintensityreadout}
  \lean{PhyXMiniProblems.ProblemPhyXMini0562.particleIntensityReadout}
  \uses{def:physics:phyx-mini-0562:phyxminiproblems-problemphyxmini0562-particleintensity}
  Read incident intensity in particles per selected square-length per selected time unit
\end{definition}
\begin{proof}
  This is the declared model definition of particle Intensity Readout.
\end{proof}
\begin{definition}[particle Rate Readout]
  \label{def:physics:phyx-mini-0562:phyxminiproblems-problemphyxmini0562-particleratereadout}
  \lean{PhyXMiniProblems.ProblemPhyXMini0562.particleRateReadout}
  \uses{def:physics:phyx-mini-0562:phyxminiproblems-problemphyxmini0562-particlerate}
  Read a particle rate in particles per selected time unit
\end{definition}
\begin{proof}
  This is the declared model definition of particle Rate Readout.
\end{proof}
\begin{definition}[length In Meters]
  \label{def:physics:phyx-mini-0562:phyxminiproblems-problemphyxmini0562-lengthinmeters}
  \lean{PhyXMiniProblems.ProblemPhyXMini0562.lengthInMeters}
  \uses{def:physics:phyx-mini-0562:phyxminiproblems-problemphyxmini0562-lengthmagnitude, def:physics:phyx-mini-0562:phyxminiproblems-problemphyxmini0562-lengthreadout}
  Meter readout used for the figure lengths R and b
\end{definition}
\begin{proof}
  This is the declared model definition of length In Meters.
\end{proof}
\begin{definition}[area In Square Meters]
  \label{def:physics:phyx-mini-0562:phyxminiproblems-problemphyxmini0562-areainsquaremeters}
  \lean{PhyXMiniProblems.ProblemPhyXMini0562.areaInSquareMeters}
  \uses{def:physics:phyx-mini-0562:phyxminiproblems-problemphyxmini0562-areamagnitude, def:physics:phyx-mini-0562:phyxminiproblems-problemphyxmini0562-areareadout}
  Square-meter readout of an effective scattering cross-section
\end{definition}
\begin{proof}
  This is the declared model definition of area In Square Meters.
\end{proof}
\begin{definition}[intensity In Particles Per Second Square Meter]
  \label{def:physics:phyx-mini-0562:phyxminiproblems-problemphyxmini0562-intensityinparticlespersecondsquaremeter}
  \lean{PhyXMiniProblems.ProblemPhyXMini0562.intensityInParticlesPerSecondSquareMeter}
  \uses{def:physics:phyx-mini-0562:phyxminiproblems-problemphyxmini0562-particleintensity, def:physics:phyx-mini-0562:phyxminiproblems-problemphyxmini0562-particleintensityreadout}
  The scalar readout of the stated incoming intensity I₀
\end{definition}
\begin{proof}
  This is the declared model definition of intensity In Particles Per Second Square Meter.
\end{proof}
\begin{definition}[rate In Particles Per Second]
  \label{def:physics:phyx-mini-0562:phyxminiproblems-problemphyxmini0562-rateinparticlespersecond}
  \lean{PhyXMiniProblems.ProblemPhyXMini0562.rateInParticlesPerSecond}
  \uses{def:physics:phyx-mini-0562:phyxminiproblems-problemphyxmini0562-particlerate, def:physics:phyx-mini-0562:phyxminiproblems-problemphyxmini0562-particleratereadout}
  The requested number of scattered particles per second
\end{definition}
\begin{proof}
  This is the declared model definition of rate In Particles Per Second.
\end{proof}
\begin{definition}[Figure Ray]
  \label{def:physics:phyx-mini-0562:phyxminiproblems-problemphyxmini0562-figureray}
  \lean{PhyXMiniProblems.ProblemPhyXMini0562.FigureRay}
  The three directed lines drawn at the collision point
\end{definition}
\begin{proof}
  This is the declared model definition of Figure Ray.
\end{proof}
\begin{definition}[Figure Length Label]
  \label{def:physics:phyx-mini-0562:phyxminiproblems-problemphyxmini0562-figurelengthlabel}
  \lean{PhyXMiniProblems.ProblemPhyXMini0562.FigureLengthLabel}
  The two physical lengths labeled in the supplied image
\end{definition}
\begin{proof}
  This is the declared model definition of Figure Length Label.
\end{proof}
\begin{definition}[Figure Angle Label]
  \label{def:physics:phyx-mini-0562:phyxminiproblems-problemphyxmini0562-figureanglelabel}
  \lean{PhyXMiniProblems.ProblemPhyXMini0562.FigureAngleLabel}
  The two copies of β and the scattering angle θ in the image
\end{definition}
\begin{proof}
  This is the declared model definition of Figure Angle Label.
\end{proof}
\begin{definition}[Hard Sphere Scattering Figure]
  \label{def:physics:phyx-mini-0562:phyxminiproblems-problemphyxmini0562-hardspherescatteringfigure}
  \lean{PhyXMiniProblems.ProblemPhyXMini0562.HardSphereScatteringFigure}
  \uses{def:physics:phyx-mini-0562:phyxminiproblems-problemphyxmini0562-lengthmagnitude, def:physics:phyx-mini-0562:phyxminiproblems-problemphyxmini0562-figureray, def:physics:phyx-mini-0562:phyxminiproblems-problemphyxmini0562-figurelengthlabel, def:physics:phyx-mini-0562:phyxminiproblems-problemphyxmini0562-figureanglelabel}
  Typed data carried by the primary bitmap. The image is schematic and has no quantitative scale; its length and angle maps state what each printed symbol denotes in the physical setup
\end{definition}
\begin{proof}
  This is the declared model definition of Hard Sphere Scattering Figure.
\end{proof}
\begin{definition}[Hard Sphere Scattering Setup]
  \label{def:physics:phyx-mini-0562:phyxminiproblems-problemphyxmini0562-hardspherescatteringsetup}
  \lean{PhyXMiniProblems.ProblemPhyXMini0562.HardSphereScatteringSetup}
  \uses{def:physics:phyx-mini-0562:phyxminiproblems-problemphyxmini0562-lengthmagnitude, def:physics:phyx-mini-0562:phyxminiproblems-problemphyxmini0562-areamagnitude, def:physics:phyx-mini-0562:phyxminiproblems-problemphyxmini0562-particleintensity, def:physics:phyx-mini-0562:phyxminiproblems-problemphyxmini0562-particlerate, def:physics:phyx-mini-0562:phyxminiproblems-problemphyxmini0562-crosssectionnormalization, def:physics:phyx-mini-0562:phyxminiproblems-problemphyxmini0562-hardspherescatteringfigure}
  All independent physical quantities in the scattering experiment. In particular, scatteredRateAboveTheta is an unknown observable: it is not defined from an answer choice or from the desired closed form
\end{definition}
\begin{proof}
  This is the declared model definition of Hard Sphere Scattering Setup.
\end{proof}
\begin{definition}[Matches Hard Sphere Scattering Scenario]
  \label{def:physics:phyx-mini-0562:phyxminiproblems-problemphyxmini0562-matcheshardspherescatteringscenario}
  \lean{PhyXMiniProblems.ProblemPhyXMini0562.MatchesHardSphereScatteringScenario}
  \uses{def:physics:phyx-mini-0562:phyxminiproblems-problemphyxmini0562-hardspherescatteringsetup}
  The qualitative physical scenario from the prose. "Negligible radius" is recorded as use of the point-particle approximation rather than as an exact claim that a physical radius equals zero
\end{definition}
\begin{proof}
  This is the declared model definition of Matches Hard Sphere Scattering Scenario.
\end{proof}
\begin{definition}[Matches Primary Figure]
  \label{def:physics:phyx-mini-0562:phyxminiproblems-problemphyxmini0562-matchesprimaryfigure}
  \lean{PhyXMiniProblems.ProblemPhyXMini0562.MatchesPrimaryFigure}
  \uses{def:physics:phyx-mini-0562:phyxminiproblems-problemphyxmini0562-hardspherescatteringsetup}
  Exact symbolic and geometric evidence read from image 562
\end{definition}
\begin{proof}
  This is the declared model definition of Matches Primary Figure.
\end{proof}
\begin{definition}[Has Physical Hard Sphere Parameters]
  \label{def:physics:phyx-mini-0562:phyxminiproblems-problemphyxmini0562-hasphysicalhardsphereparameters}
  \lean{PhyXMiniProblems.ProblemPhyXMini0562.HasPhysicalHardSphereParameters}
  \uses{def:physics:phyx-mini-0562:phyxminiproblems-problemphyxmini0562-lengthinmeters, def:physics:phyx-mini-0562:phyxminiproblems-problemphyxmini0562-intensityinparticlespersecondsquaremeter, def:physics:phyx-mini-0562:phyxminiproblems-problemphyxmini0562-hardspherescatteringsetup}
  Physical angle and magnitude ranges. Angles are measured in radians: β lies between zero and a right angle and θ between zero and π
\end{definition}
\begin{proof}
  This is the declared model definition of Has Physical Hard Sphere Parameters.
\end{proof}
\begin{definition}[Satisfies Specular Reflection Law]
  \label{def:physics:phyx-mini-0562:phyxminiproblems-problemphyxmini0562-satisfiesspecularreflectionlaw}
  \lean{PhyXMiniProblems.ProblemPhyXMini0562.SatisfiesSpecularReflectionLaw}
  \uses{def:physics:phyx-mini-0562:phyxminiproblems-problemphyxmini0562-hardspherescatteringsetup}
  Specular reflection at a smooth hard surface. The first field is the equal angle law. The second expresses the deflection between the incoming and outgoing rays before substituting that equality; it does not contain the requested particle rate
\end{definition}
\begin{proof}
  This is the declared model definition of Satisfies Specular Reflection Law.
\end{proof}
\begin{definition}[Satisfies Hard Sphere Contact Geometry]
  \label{def:physics:phyx-mini-0562:phyxminiproblems-problemphyxmini0562-satisfieshardspherecontactgeometry}
  \lean{PhyXMiniProblems.ProblemPhyXMini0562.SatisfiesHardSphereContactGeometry}
  \uses{def:physics:phyx-mini-0562:phyxminiproblems-problemphyxmini0562-lengthinmeters, def:physics:phyx-mini-0562:phyxminiproblems-problemphyxmini0562-hardspherescatteringsetup}
  Right-triangle contact geometry from the figure. The perpendicular impact parameter is the radius times the sine of the incidence angle
\end{definition}
\begin{proof}
  This is the declared model definition of Satisfies Hard Sphere Contact Geometry.
\end{proof}
\begin{definition}[Satisfies Uniform Beam Counting Law]
  \label{def:physics:phyx-mini-0562:phyxminiproblems-problemphyxmini0562-satisfiesuniformbeamcountinglaw}
  \lean{PhyXMiniProblems.ProblemPhyXMini0562.SatisfiesUniformBeamCountingLaw}
  \uses{def:physics:phyx-mini-0562:phyxminiproblems-problemphyxmini0562-areainsquaremeters, def:physics:phyx-mini-0562:phyxminiproblems-problemphyxmini0562-intensityinparticlespersecondsquaremeter, def:physics:phyx-mini-0562:phyxminiproblems-problemphyxmini0562-rateinparticlespersecond, def:physics:phyx-mini-0562:phyxminiproblems-problemphyxmini0562-hardspherescatteringsetup}
  Uniform-beam counting: scattered rate equals incoming intensity times the effective cross-section for the selected angular range. This is a generic flux-times-area law and contains no answer-choice expression
\end{definition}
\begin{proof}
  This is the declared model definition of Satisfies Uniform Beam Counting Law.
\end{proof}
\begin{definition}[Satisfies Hard Sphere Cross Section Law]
  \label{def:physics:phyx-mini-0562:phyxminiproblems-problemphyxmini0562-satisfieshardspherecrosssectionlaw}
  \lean{PhyXMiniProblems.ProblemPhyXMini0562.SatisfiesHardSphereCrossSectionLaw}
  \uses{def:physics:phyx-mini-0562:phyxminiproblems-problemphyxmini0562-crosssectionnormalizationfactor, def:physics:phyx-mini-0562:phyxminiproblems-problemphyxmini0562-lengthinmeters, def:physics:phyx-mini-0562:phyxminiproblems-problemphyxmini0562-areainsquaremeters, def:physics:phyx-mini-0562:phyxminiproblems-problemphyxmini0562-hardspherescatteringsetup}
  Hard-sphere impact-disk law for the angular region above the threshold. The law is stated before eliminating b: its area is π b² for an ordinary three-dimensional beam and b² under the convention implicit in the displayed choices. Thus this premise does not contain any final rate formula or any answer-choice angular expression
\end{definition}
\begin{proof}
  This is the declared model definition of Satisfies Hard Sphere Cross Section Law.
\end{proof}
\begin{definition}[Answer Choice]
  \label{def:physics:phyx-mini-0562:phyxminiproblems-problemphyxmini0562-answerchoice}
  \lean{PhyXMiniProblems.ProblemPhyXMini0562.AnswerChoice}
  Labels of the four answer choices in the source problem
\end{definition}
\begin{proof}
  This is the declared model definition of Answer Choice.
\end{proof}
\begin{definition}[displayed Angular Factor]
  \label{def:physics:phyx-mini-0562:phyxminiproblems-problemphyxmini0562-displayedangularfactor}
  \lean{PhyXMiniProblems.ProblemPhyXMini0562.displayedAngularFactor}
  \uses{def:physics:phyx-mini-0562:phyxminiproblems-problemphyxmini0562-answerchoice}
  The dimensionless angular factor printed in each answer choice
\end{definition}
\begin{proof}
  This is the declared model definition of displayed Angular Factor.
\end{proof}
\begin{definition}[displayed Rate Expression]
  \label{def:physics:phyx-mini-0562:phyxminiproblems-problemphyxmini0562-displayedrateexpression}
  \lean{PhyXMiniProblems.ProblemPhyXMini0562.displayedRateExpression}
  \uses{def:physics:phyx-mini-0562:phyxminiproblems-problemphyxmini0562-lengthinmeters, def:physics:phyx-mini-0562:phyxminiproblems-problemphyxmini0562-intensityinparticlespersecondsquaremeter, def:physics:phyx-mini-0562:phyxminiproblems-problemphyxmini0562-hardspherescatteringsetup, def:physics:phyx-mini-0562:phyxminiproblems-problemphyxmini0562-answerchoice, def:physics:phyx-mini-0562:phyxminiproblems-problemphyxmini0562-displayedangularfactor}
  The particle-rate readout printed beside a selected answer label
\end{definition}
\begin{proof}
  This is the declared model definition of displayed Rate Expression.
\end{proof}
\begin{definition}[recorded Dataset Answer]
  \label{def:physics:phyx-mini-0562:phyxminiproblems-problemphyxmini0562-recordeddatasetanswer}
  \lean{PhyXMiniProblems.ProblemPhyXMini0562.recordedDatasetAnswer}
  \uses{def:physics:phyx-mini-0562:phyxminiproblems-problemphyxmini0562-answerchoice}
  The answer label recorded by the dataset; it is not a theorem premise
\end{definition}
\begin{proof}
  This is the declared model definition of recorded Dataset Answer.
\end{proof}
\begin{definition}[Matches Answer Choice]
  \label{def:physics:phyx-mini-0562:phyxminiproblems-problemphyxmini0562-matchesanswerchoice}
  \lean{PhyXMiniProblems.ProblemPhyXMini0562.MatchesAnswerChoice}
  \uses{def:physics:phyx-mini-0562:phyxminiproblems-problemphyxmini0562-rateinparticlespersecond, def:physics:phyx-mini-0562:phyxminiproblems-problemphyxmini0562-hardspherescatteringsetup, def:physics:phyx-mini-0562:phyxminiproblems-problemphyxmini0562-answerchoice, def:physics:phyx-mini-0562:phyxminiproblems-problemphyxmini0562-displayedrateexpression}
  A displayed answer agrees with the modeled scattered particle rate
\end{definition}
\begin{proof}
  This is the declared model definition of Matches Answer Choice.
\end{proof}
\begin{lemma}[scattering angle eq pi sub two beta]
  \label{lem:physics:phyx-mini-0562:phyxminiproblems-problemphyxmini0562-scattering-angle-eq-pi-sub-two-beta}
  \lean{PhyXMiniProblems.ProblemPhyXMini0562.scattering_angle_eq_pi_sub_two_beta}
  \uses{def:physics:phyx-mini-0562:phyxminiproblems-problemphyxmini0562-hardspherescatteringsetup, def:physics:phyx-mini-0562:phyxminiproblems-problemphyxmini0562-satisfiesspecularreflectionlaw}
  The specular-reflection law gives the angle relation stated in the problem, θ = π - 2β. This is a derived geometric relation, not a rate formula
\end{lemma}
\begin{proof}
  The conclusion follows from the governing relations and figure readouts named in the statement of scattering angle eq pi sub two beta.
\end{proof}
\begin{lemma}[impact parameter eq radius mul cos half scattering angle]
  \label{lem:physics:phyx-mini-0562:phyxminiproblems-problemphyxmini0562-impact-parameter-eq-radius-mul-cos-half-scattering-angle}
  \lean{PhyXMiniProblems.ProblemPhyXMini0562.impact_parameter_eq_radius_mul_cos_half_scattering_angle}
  \uses{def:physics:phyx-mini-0562:phyxminiproblems-problemphyxmini0562-lengthinmeters, def:physics:phyx-mini-0562:phyxminiproblems-problemphyxmini0562-hardspherescatteringsetup, def:physics:phyx-mini-0562:phyxminiproblems-problemphyxmini0562-satisfiesspecularreflectionlaw, def:physics:phyx-mini-0562:phyxminiproblems-problemphyxmini0562-satisfieshardspherecontactgeometry}
  Combining the scattering-angle relation with the contact triangle gives b = R cos(θ/2). The result still does not determine a particle rate
\end{lemma}
\begin{proof}
  The conclusion follows from the governing relations and figure readouts named in the statement of impact parameter eq radius mul cos half scattering angle.
\end{proof}
\begin{lemma}[effective cross section above theta]
  \label{lem:physics:phyx-mini-0562:phyxminiproblems-problemphyxmini0562-effective-cross-section-above-theta}
  \lean{PhyXMiniProblems.ProblemPhyXMini0562.effective_cross_section_above_theta}
  \uses{def:physics:phyx-mini-0562:phyxminiproblems-problemphyxmini0562-crosssectionnormalizationfactor, def:physics:phyx-mini-0562:phyxminiproblems-problemphyxmini0562-lengthinmeters, def:physics:phyx-mini-0562:phyxminiproblems-problemphyxmini0562-areainsquaremeters, def:physics:phyx-mini-0562:phyxminiproblems-problemphyxmini0562-hardspherescatteringsetup, def:physics:phyx-mini-0562:phyxminiproblems-problemphyxmini0562-satisfiesspecularreflectionlaw, def:physics:phyx-mini-0562:phyxminiproblems-problemphyxmini0562-satisfieshardspherecontactgeometry, def:physics:phyx-mini-0562:phyxminiproblems-problemphyxmini0562-satisfieshardspherecrosssectionlaw}
  Eliminating the impact parameter gives the convention-independent shape of the effective cross-section. Its normalization remains explicit: π for a full impact disk and 1 for the source convention
\end{lemma}
\begin{proof}
  The conclusion follows from the governing relations and figure readouts named in the statement of effective cross section above theta.
\end{proof}
\begin{theorem}[number scattered per second above theta full impact disk]
  \label{thm:physics:phyx-mini-0562:phyxminiproblems-problemphyxmini0562-number-scattered-per-second-above-theta-full-impact-disk}
  \lean{PhyXMiniProblems.ProblemPhyXMini0562.number_scattered_per_second_above_theta_full_impact_disk}
  \uses{def:physics:phyx-mini-0562:phyxminiproblems-problemphyxmini0562-lengthinmeters, def:physics:phyx-mini-0562:phyxminiproblems-problemphyxmini0562-intensityinparticlespersecondsquaremeter, def:physics:phyx-mini-0562:phyxminiproblems-problemphyxmini0562-rateinparticlespersecond, def:physics:phyx-mini-0562:phyxminiproblems-problemphyxmini0562-hardspherescatteringsetup, def:physics:phyx-mini-0562:phyxminiproblems-problemphyxmini0562-matcheshardspherescatteringscenario, def:physics:phyx-mini-0562:phyxminiproblems-problemphyxmini0562-matchesprimaryfigure, def:physics:phyx-mini-0562:phyxminiproblems-problemphyxmini0562-hasphysicalhardsphereparameters, def:physics:phyx-mini-0562:phyxminiproblems-problemphyxmini0562-satisfiesspecularreflectionlaw, def:physics:phyx-mini-0562:phyxminiproblems-problemphyxmini0562-satisfieshardspherecontactgeometry, def:physics:phyx-mini-0562:phyxminiproblems-problemphyxmini0562-satisfiesuniformbeamcountinglaw, def:physics:phyx-mini-0562:phyxminiproblems-problemphyxmini0562-satisfieshardspherecrosssectionlaw}
  With ordinary particles-per-second-per-square-metre intensity and the full three-dimensional impact disk, the rate contains the geometrically required factor π
\end{theorem}
\begin{proof}
  The conclusion follows from the governing relations and figure readouts named in the statement of number scattered per second above theta full impact disk.
\end{proof}
\begin{theorem}[number scattered per second above theta source reduced]
  \label{thm:physics:phyx-mini-0562:phyxminiproblems-problemphyxmini0562-number-scattered-per-second-above-theta-source-reduced}
  \lean{PhyXMiniProblems.ProblemPhyXMini0562.number_scattered_per_second_above_theta_source_reduced}
  \uses{def:physics:phyx-mini-0562:phyxminiproblems-problemphyxmini0562-lengthinmeters, def:physics:phyx-mini-0562:phyxminiproblems-problemphyxmini0562-intensityinparticlespersecondsquaremeter, def:physics:phyx-mini-0562:phyxminiproblems-problemphyxmini0562-rateinparticlespersecond, def:physics:phyx-mini-0562:phyxminiproblems-problemphyxmini0562-hardspherescatteringsetup, def:physics:phyx-mini-0562:phyxminiproblems-problemphyxmini0562-matcheshardspherescatteringscenario, def:physics:phyx-mini-0562:phyxminiproblems-problemphyxmini0562-matchesprimaryfigure, def:physics:phyx-mini-0562:phyxminiproblems-problemphyxmini0562-hasphysicalhardsphereparameters, def:physics:phyx-mini-0562:phyxminiproblems-problemphyxmini0562-satisfiesspecularreflectionlaw, def:physics:phyx-mini-0562:phyxminiproblems-problemphyxmini0562-satisfieshardspherecontactgeometry, def:physics:phyx-mini-0562:phyxminiproblems-problemphyxmini0562-satisfiesuniformbeamcountinglaw, def:physics:phyx-mini-0562:phyxminiproblems-problemphyxmini0562-satisfieshardspherecrosssectionlaw}
  Under the reduced convention used by all four displayed options, the same model gives the source's recorded expression without π
\end{theorem}
\begin{proof}
  The conclusion follows from the governing relations and figure readouts named in the statement of number scattered per second above theta source reduced.
\end{proof}
\begin{theorem}[answer choice D matches under source normalization]
  \label{thm:physics:phyx-mini-0562:phyxminiproblems-problemphyxmini0562-answer-choice-d-matches-under-source-normalization}
  \lean{PhyXMiniProblems.ProblemPhyXMini0562.answer_choice_D_matches_under_source_normalization}
  \uses{def:physics:phyx-mini-0562:phyxminiproblems-problemphyxmini0562-hardspherescatteringsetup, def:physics:phyx-mini-0562:phyxminiproblems-problemphyxmini0562-matcheshardspherescatteringscenario, def:physics:phyx-mini-0562:phyxminiproblems-problemphyxmini0562-matchesprimaryfigure, def:physics:phyx-mini-0562:phyxminiproblems-problemphyxmini0562-hasphysicalhardsphereparameters, def:physics:phyx-mini-0562:phyxminiproblems-problemphyxmini0562-satisfiesspecularreflectionlaw, def:physics:phyx-mini-0562:phyxminiproblems-problemphyxmini0562-satisfieshardspherecontactgeometry, def:physics:phyx-mini-0562:phyxminiproblems-problemphyxmini0562-satisfiesuniformbeamcountinglaw, def:physics:phyx-mini-0562:phyxminiproblems-problemphyxmini0562-satisfieshardspherecrosssectionlaw, def:physics:phyx-mini-0562:phyxminiproblems-problemphyxmini0562-matchesanswerchoice}
  Consequently, the modeled source-reduced rate agrees with displayed choice D. The normalization hypothesis is convention data, not the desired rate equality
\end{theorem}
\begin{proof}
  The conclusion follows from the governing relations and figure readouts named in the statement of answer choice D matches under source normalization.
\end{proof}
% --- Archon declaration topology end ---
% --- Archon physics formalization source end ---
